% !TeX encoding = UTF-8
% !TeX spellcheck = de_DE
\pdfminorversion=4

%\newcommand{\todo}[1]{\textbf{\color{red}{TODO: #1}}}

\documentclass[german]{beamer} % evtl. FabLab Anpassungen
\usepackage[utf8]{inputenc}
\usepackage[german]{babel}
\usepackage{color}
\usepackage{amsmath, amssymb, amstext, amsfonts, mathrsfs}
\usepackage{hyperref}
\usepackage{scrhack}
\usepackage{multicol}

\newcommand{\images}{vektorzeichnen/}
\newcommand{\ccbysa}{\href{https://creativecommons.org/licenses/by-sa/3.0/deed.de}{\includegraphics[height=0.11in]{licenses/cc-by-sa.png}}}
\newcommand{\artlibre}{\href{http://artlibre.org/licence/lal/de/}{\includegraphics[height=0.11in]{licenses/art_libre.png}}}

\beamertemplatenavigationsymbolsempty

\addtobeamertemplate{footline}{
    \usebeamerfont{footline}
    \usebeamercolor[fg]{footline}
    \hspace{1em}%
    \insertframenumber/\inserttotalframenumber
    \hspace{1em}%
    \insertsectionhead
}{}

%-----------------------------------------------------------------------------------------------------------------

\title{Vektorzeichnen mit Inkscape}
\author{Sebastian Endres}
\logo{\includegraphics[width=0.3\linewidth]{logo/logo.pdf}}

\begin{document}

%-----------------------------------------------------------------------------------------------------------------

\maketitle

%-----------------------------------------------------------------------------------------------------------------

\frame{\tableofcontents}

%-----------------------------------------------------------------------------------------------------------------

\section{Grundlagen Vektorgrafik}
\subsection{Vergleich Pixelgrafik - Vektorgrafik}
\begin{frame}[t]
  \frametitle{Was ist eine Vektorzeichnung?}
  
  \begin{columns}[t]
  	\begin{column}{0.5\textwidth}
  		\textbf{Pixelgrafik}
  		\begin{itemize}
  			\item Besteht aus einzelnen Pixeln (Farbpunkte)
  			\item verpixelt beim Vergrößern
  			\item Bsp.: \texttt{.png, .jp(e)g, .bmp, .ico, .gif}
  		\end{itemize}
  		\begin{figure}
  		   \centering
  			\includegraphics[width=0.54\textwidth]{\images/gnu_bitter.png}
  			\caption{verpixelter GNU}
  			\label{fig:gnu_bitter}
  		\end{figure}
  	\end{column}
  	
  	\begin{column}{0.5\textwidth}
  	\textbf{Vektorgrafik}
  		\begin{itemize}
  			\item Besteht aus geometrischen Objekten \\(Linien, Flächen, etc.)
  			\item Objekte bestehen aus Pfaden
  			\item Bsp.: \texttt{.svg}
  		\end{itemize}
  		\begin{figure}
  			\centering
  			\includegraphics[width=0.4\textwidth]{\images/gnu.pdf}
  			\caption{vektorisierter GNU}
  			\label{fig:gnu_vektor}
  		\end{figure}
  	\end{column}
  \end{columns}
  
 \end{frame}

%-----------------------------------------------------------------------------------------------------------------

\subsection{Einsatzgebiete für Vektorgrafiken?}
\begin{frame}
  \frametitle{Warum Vektorgrafik?}
  
  \begin{itemize}
	  \item Aus Vektorpfaden kann der Weg für Maschinen (Lasercutter / Schneideplotter) berechnet werden
	  		\begin{minipage}{\linewidth}
	  		  \begin{figure}
	  		  	\centering
	  		    \includegraphics[width=0.3\textwidth]{\images/path.pdf}
	  		    \caption{Pfad}
	  		    \label{fig:path}
	  		  \end{figure}
	  		\end{minipage}
	  \item Vektorgrafiken sind beliebig skalierbar
  \end{itemize}

  
 \end{frame}

%-----------------------------------------------------------------------------------------------------------------

\subsection{Programme - Übersicht}
\begin{frame}
  \frametitle{Welche Programme gibt es?}
  
  \begin{itemize}
  	\item Inkscape
  	\item Adobe Illustrator (proprietär, kostenpflichtig)
  	\item Corel Draw (proprietär, kostenpflichtig)
  \end{itemize}
  
 \end{frame}

%-----------------------------------------------------------------------------------------------------------------

\section{Inkscape}
\subsection{Grundlagen}
\begin{frame}
  \frametitle{Inkscape}
  
  \begin{itemize}
  	\item Kostenlos
  	\item open source
  	\item Ideal für das Erstellen von SVG-Dateien
  	\item Verfügbar für Linux, Windows, Mac
  	\item \href{https://inkscape.org/download/}{Download Seite: \texttt{inkscape.org/download}}
  \end{itemize}

  \begin{figure}
  	\centering
    \includegraphics[width=0.3\textwidth]{\images/inkscape_logo_300x300.png}
    \caption{Inkscape Logo}
    \label{fig:inkscape_logo}
  \end{figure}  
  
 \end{frame}

%-----------------------------------------------------------------------------------------------------------------

\subsubsection{Übersicht}
\begin{frame}
  \frametitle{Übersicht}

  \begin{figure}
  	\centering
    \includegraphics[width=\textwidth]{\images/scr_shot/start.png}
    %\caption{Start}
    \label{fig:inkscape_scr_shot_start}
  \end{figure}
  
 \end{frame}

 
%-----------------------------------------------------------------------------------------------------------------

\subsubsection{Objekte zeichnen}
\begin{frame}
  \frametitle{Objekt zeichnen}

  \begin{figure}
  	\centering
    \includegraphics[width=\textwidth]{\images/scr_shot/draw_object.png}
    %\caption{Object zeichnen}
    \label{fig:inkscape_scr_shot_draw_object}
  \end{figure}
  
 \end{frame}
 
%-----------------------------------------------------------------------------------------------------------------

\subsubsection{Farbe, Dicke setzen}
\begin{frame}
  \frametitle{Farbe, Dicke setzen}

  \begin{figure}
  	\centering
    \includegraphics[width=0.7\textwidth]{\images/scr_shot/color.png}
    \caption{Farbe setzen}
    \label{fig:inkscape_scr_shot_color}
  \end{figure}
      
  \textbf{Farb- und Linien Fenster:}
  \begin{itemize}
  	\item \texttt{Objekt $\rightarrow$ Füllung und Kontur}
  	\item \texttt{Strg + Umschalt + F}
	\item Doppelklick auf \colorbox{orange}{Farbbereich}
  \end{itemize} 
  
 \end{frame}
 
%-----------------------------------------------------------------------------------------------------------------

\subsubsection{Pfad manuell bearbeiten}
\begin{frame}[t]
  \frametitle{Pfad manuell bearbeiten}
  
  \begin{columns}[t]
  	\begin{column}{0.4\textwidth}
  		\textbf{Mäuse}
  		\begin{figure}
  		   \centering
  			\includegraphics[width=0.3\textwidth]{\images/scr_shot/mouses.png}
  			\label{fig:inkscape_scr_shot_mouses}
  		\end{figure}
  		\begin{enumerate}
  			\item Objekt Maus für das Bearbeiten von Objekten
  			\item Pfad Maus für das Bearbeiten von Pfaden anhand ihrer Knoten
  		\end{enumerate}
  	\end{column}
  	
  	\begin{column}{0.6\textwidth}
  	\textbf{Pfad bearbeiten}
  		\begin{figure}
  			\centering
  			\includegraphics[width=0.9\textwidth]{\images/scr_shot/edit_path.png}
  			\label{fig:inkscape_scr_shot_edit_path}
  		\end{figure}
  	\end{column}
  \end{columns}
  
 \end{frame}
 
%-----------------------------------------------------------------------------------------------------------------

\subsubsection{Pfad Optionen}
\begin{frame}[t]
  \frametitle{Pfad Optionen}
  
  \begin{columns}[t]
  	\begin{column}{0.3\textwidth}
  		\begin{figure}
  		   \centering
  			\includegraphics[width=\textwidth]{\images/scr_shot/path_options.png}
  			\label{fig:inkscape_scr_shot_path_options}
  		\end{figure}
  	\end{column}
  	
  	\begin{column}{0.7\textwidth}
  		\begin{enumerate}
  			\item Mehrere Pfad-Objekte können zu einem Objekt vereint werden
  			\item Sie können einfach kombiniert werden (\texttt{Strg + K}), dann sind aber die einzelnen Linien nicht verbunden sondern nur \grqq zusammengehörig\grqq
  			\item Sie können \textbf{verschnitten} werden (verschiedene \grqq Regeln\grqq)
  		\end{enumerate}
  		
  		\vspace{1em}
  		\alert{\textbf{Beachte:} Gruppieren verbindet Objekte, aber vereint sie nicht.}
  		
  	\end{column}
  \end{columns}
  
 \end{frame}
 
%-----------------------------------------------------------------------------------------------------------------

\subsubsection{Bilder vektorisieren}
\begin{frame}
	\frametitle{Bilder Vektorisieren}
  	\centering
	\textbf{Pixelbilder können in Vektorbilder umgewandelt werden}
  
	\begin{figure}
  		\centering
    	\includegraphics[height=0.3\textwidth]{\images/scr_shot/trace_bitmap.png}
    	\caption{Inkscape Screenshot - Bitmap vektorisieren}
    	\label{fig:inkscape_scr_shot_trace_bitmap}
  	\end{figure}
  	
  	\begin{itemize}
  		\item \texttt{Pfad $\rightarrow$ Bitmap vektorisieren}
  		\item \texttt{Alt + Shift + B}
  	\end{itemize}
  
 \end{frame}

%-----------------------------------------------------------------------------------------------------------------

\subsection{Beachte beim Zeichnen}
\begin{frame}
  \frametitle{Beachte beim Zeichnen}
  
  \begin{itemize}
  	\item Dünne Linien verwenden, um genauer zeichnen zu können
  	\item Objekte ausrichten und einrasten lassen
  		\begin{itemize}
  			\item Verwende die Einrastoptionen am rechten Rand
  			\item Drücke \texttt{Strg} beim Verschieben, um die Richtungen zu sperren
  			\item Richte die Objekte aus (\texttt{Objekt $\rightarrow$ Ausrichten und anordnen})
  		\end{itemize}
	\item Zusammengehörige Objekte verbinden / vereinen, nicht nur gruppieren
  \end{itemize} 
  
 \end{frame}

%-----------------------------------------------------------------------------------------------------------------












 
%-----------------------------------------------------------------------------------------------------------------

\section{Obligatorische Schlussfolie}
\begin{frame}
  \frametitle{Obligatorische Schlussfolie}
  
  \begin{block}{Vielen Dank für die Aufmerksamkeit}\end{block}
  
  \textbf{Quellen}
  \begin{itemize}
  	\item \textbf{Inkscape Official Logo (Raster, 300x300px)}\\ von \textit{Inkscape User}: \href{https://inkscape.org/en/gallery/item/1088/}{\texttt{inkscape.org}} (Lizenz: \ccbysa) \par
  	        \begin{minipage}{\linewidth}
  	            \centering
  	            \includegraphics[width=0.5in]{\images/inkscape_logo_300x300.png}
  	        \end{minipage}
  			

  	\item \textbf{Heckert GNU white}\\ von \textit{Anomie}: \href{https://en.wikipedia.org/wiki/File:Heckert_GNU_white.svg}{\texttt{wikipedia.org}} (Lizenz: \artlibre) \par
  	        \begin{minipage}{\linewidth}
  	            \centering
  	            \includegraphics[width=0.5in]{\images/gnu.png}
  	        \end{minipage}
  			
  	\item \textbf{Eigene Bilder} (Lizenz: \ccbysa)
  \end{itemize}
  
  \begin{block}{Diese Präsentation ist verfügbar unter einer \ccbysa $\ $ 3.0 Lizenz.}\end{block}
  
 \end{frame}

%-----------------------------------------------------------------------------------------------------------------

\end{document}
